% Living discrepancy ledger for the four primary DKPS theorem targets.
% Update this file whenever a public theorem changes its assumptions,
% conclusion, estimator, asymptotic regime, or relationship to the source paper.

\documentclass[11pt]{article}

\usepackage[T1]{fontenc}
\usepackage{lmodern}
\usepackage{microtype}
\usepackage[margin=1.0in]{geometry}
\usepackage{amsmath,amssymb,mathtools}
\usepackage{booktabs,longtable,array}
\usepackage[colorlinks=true,linkcolor=black,urlcolor=blue,citecolor=blue]{hyperref}

\newcommand{\Lean}{\textsc{Lean}}
\newcommand{\DKPS}{\textsc{DKPS}}
\newcommand{\whp}{\text{with high probability}}
\newcommand{\norm}[1]{\left\lVert #1\right\rVert}
\newcommand{\code}[1]{\texttt{#1}}
\setlength{\tabcolsep}{4pt}
\setlength{\emergencystretch}{2em}

\title{The Four Main DKPS Papers:\\Differences Between the Literature and the Lean Formalization}
\author{Jonathan Crall \and GPT-5.6 Thinking}
\date{\today}

\begin{document}
\maketitle

\begin{abstract}
This is a living audit of the four primary paper targets in the
\DKPS{} formalization repository.  Its purpose is not to restate every theorem,
but to record every material difference between the published mathematical
presentation and the theorem that is actually proved in \Lean.  Differences are
classified as repaired statements, surfaced assumptions, strengthened or
weakened conclusions, modelling choices, and still-open composition gaps.  The
source of truth remains the theorem signatures and proofs; this document is the
human-readable discrepancy ledger that should change with them.
\end{abstract}

\section{Scope and maintenance rule}

The four direct targets are:
\begin{enumerate}
\item Acharyya, Trosset, Priebe, and Helm (2024), consistent estimation of
      generative-model representations in the data kernel perspective space;
\item Acharyya, Agterberg, Park, and Priebe (2025), concentration bounds for
      response-based vector embeddings;
\item Helm, Acharyya, Park, Duderstadt, and Priebe (2025), statistical inference
      on black-box generative models in DKPS;
\item Helm, Johnson, and Priebe (2026), query-efficient model evaluation using
      cached responses (Quench).
\end{enumerate}

A row belongs in this document when any of the following changes:
\begin{itemize}
\item an assumption present in \Lean{} is absent or only implicit in the paper;
\item a paper statement is false or underspecified without a repair;
\item the formalization proves a stronger or weaker conclusion;
\item the formalization uses a different estimator or MDS variant;
\item a probability, measurability, compactness, or selection issue is made
      explicit;
\item a formal theorem covers only a special case of the paper, or extends it.
\end{itemize}

Compatibility theorems may retain older, heavier interfaces.  The preferred
public capstone is the theorem with the smallest honest hypothesis load.  A
proof refactor that merely shortens an argument does not change this ledger; a
refactor that removes a caller-visible condition does.

\section{Acharyya et al. (2024): asymptotic DKPS and raw-stress MDS}

\begin{longtable}{>{\raggedright\arraybackslash}p{0.18\textwidth}
                  >{\raggedright\arraybackslash}p{0.25\textwidth}
                  >{\raggedright\arraybackslash}p{0.36\textwidth}
                  >{\raggedright\arraybackslash}p{0.13\textwidth}}
\toprule
Topic & Literature & Formalization & Classification \\
\midrule
\endhead
Choice of population minimizer & The fixed-configuration theorem is phrased
up to affine ambiguity, with a subsequence selecting a compatible limit. & The
unconditional theorem is set-valued: estimated minimizers approach the set of
population pair-distance profiles.  A literal fixed-configuration theorem
requires the explicit predicate \code{UniquePairProfile}. & Repaired statement
\\
\addlinespace
Subsequence versus full sequence & The paper uses subsequences and a diagonal
argument. & Once uniqueness of the pair-distance profile is assumed, the
formalization obtains full-sequence convergence; the diagonal argument is no
longer needed. & Stronger conclusion under explicit uniqueness \\
\addlinespace
Measurable selection & A measurable choice of MDS minimizer is left implicit. &
The proof uses a deterministic modulus of continuity and event inclusion, so no
measurable selector is required. & Assumption removed \\
\addlinespace
Configuration error & Informal statements identify configurations only up to
affine transformation. & \code{ConfigError} is direct coordinate error and is
therefore stronger than affine-quotiented closeness.  The repaired consistency
layer primarily compares pairwise-distance profiles where that is the invariant
quantity. & Modelling distinction \\
\addlinespace
Moment assumptions & The sampling discussion is iid. & The variance algebra is
proved under pairwise independence, which is sufficient.  Measurability,
integrability, and \(L^2\) hypotheses are explicit. & Generalization plus
surfaced technical assumptions \\
\addlinespace
Distribution-over-model conclusion & The paper also discusses an
\(L^p\)-over-model-distribution form. & That form is not yet formalized in the
2024 library. & Scope gap \\
\bottomrule
\end{longtable}

The important repair is conceptual: nonuniqueness of raw-stress MDS is treated
as a property of the theorem, not hidden inside a choice function.  The
formalization's strongest unconditional result is therefore about the minimizer
set or its invariant pair-distance profile.  A fixed representative is available
only after a uniqueness condition makes it mathematically meaningful.


\paragraph{Role in the Perfect Quench scaffold.}
The new raw-response track reuses the 2024 matrix-valued sample-mean
second-moment theorem rather than postulating response means.  It adds an
independent product probability space for reference and response randomness,
then lifts the existing replicate-average calculation through random reference
selection.  This is composition work beyond the paper, not a change to the
2024 consistency theorem itself.

\section{Acharyya et al. (2025): finite-sample response embeddings}

\begin{longtable}{>{\raggedright\arraybackslash}p{0.18\textwidth}
                  >{\raggedright\arraybackslash}p{0.25\textwidth}
                  >{\raggedright\arraybackslash}p{0.36\textwidth}
                  >{\raggedright\arraybackslash}p{0.13\textwidth}}
\toprule
Topic & Literature & Formalization & Classification \\
\midrule
\endhead
Population spectral assumptions & Rank \(d\), positive spectrum, and upper
spectral control are stated at paper level. & Positive semidefiniteness, rank,
and eigenvalue floor/cap are represented as exact matrix predicates.  In the
preferred Quench-facing capstones, a population Gram witness now derives
positive semidefiniteness and rank automatically. & Surfaced and partially
discharged \\
\addlinespace
Smallness regime & The paper packages the regime through asymptotic growth such
as sufficiently many responses relative to the model count. & The finite
spectral theorem exposes the exact entrywise, eigengap, and polar-factor
smallness inequalities consumed by the proof.  \code{GrowingConfigControl}
packages their eventual validity for varying matrix size. & Quantifiers made
explicit \\
\addlinespace
Alignment choice & The paper chooses an optimal orthogonal alignment
classically. & The formalization separates the existential alignment predicate
from a noncomputable selected alignment, so probability arguments do not depend
on measurability of a data-dependent choice. & Selection issue resolved \\
\addlinespace
Rates & The paper presents sharper asymptotic bookkeeping. &
\code{cmdsEntrywiseRate}, \code{configBound}, and the composed end-to-end rate
are valid but intentionally loose.  No claim of optimal constants or exponents
is made. & Weaker quantitative sharpness \\
\addlinespace
Growing target-augmented CMDS & The paper's theorem is presented for a fixed
finite model collection. & The formalization additionally supports a growing
\(n+1\) batch containing the \(n\) sampled references and the current target,
and proves pairwise-distance control without a globally aligned coordinate
map. & Extension beyond paper \\
\addlinespace
Response boundedness & Bounded dissimilarities are used in the concentration
pipeline. & The preferred response bridge needs only a population response-norm
envelope.  On the response-mean event, the random sample norm and both sample
and population dissimilarity bounds are derived internally. & Hypotheses
reduced \\
\addlinespace
Davis--Kahan implementation & The paper invokes spectral perturbation theory. &
The production proof currently uses a proved bespoke cross-energy and polar
factor chain.  Newer general Davis--Kahan results are intentionally deferred
unless they remove a caller-visible condition, such as the polar smallness
hypothesis. & Deliberate implementation choice \\
\bottomrule
\end{longtable}

The genuine remaining paper-facing assumption is a population covariance
floor; it is an identifiability condition for the spectral estimator rather than
a proof-engineering artifact.  The formalization now proves the first
deterministic reduction in that chain: entrywise closeness to the population
covariance preserves at least half of its quadratic-form floor.  This theorem
requires no separate compactness, measurability, or integrability premise;
integrability of the required coordinate products is recovered from the
positive nondegeneracy certificate by polarization.  The formalization now
also identifies the sample-centered empirical covariance quadratic form with
the average squared projections of the centered reference cloud, so the event
membership directly yields the reference-scatter quadratic floor.  The
rectangular scatter--Gram spectral bridge, the target-augmentation floor, and
the assembly of the measurable high-probability spectral certificate are now
proved as well: the quadratic floor transfers to the first $d$ sorted
eigenvalues of the centered reference Gram matrix through the general
nonzero-spectrum bridge between $T^*T$ and $TT^*$, augmentation preserves the
floor through the exact add-one centered-scatter identity, and the certificate
event is the covariance event intersected with the deterministic dimension
gate $d \le n$.


\paragraph{Role in the Perfect Quench scaffold.}
The completion program replaces stagewise global spectral assumptions by a
high-probability certificate derived from a population covariance floor and iid
reference sampling.  The deterministic covariance perturbation seam is now
closed in \code{DkpsQuench.Perfect.CovarianceFloor}, and the sample-centered
quadratic identity is closed in \code{DkpsQuench.Perfect.CenteredCovariance}.
The rectangular scatter/Gram eigenvalue transfers are closed in
\code{DkpsQuench.Perfect.GramSpectrumBridge} on top of the new proved
\code{ForMathlib} rectangular spectral layer, and the spectral-event assembly
\code{exists\_growingSpectralSubevents\_of\_compact\_iid\_nondegenerate} is
proved in \code{DkpsQuench.Perfect.SpectralRegularity}.  The explicit response
schedule is already proved.  These changes reduce the assumptions of the
Quench composition without altering the proved 2025 concentration statements.

\section{Helm et al. (2025): statistical inference transfer}

\begin{longtable}{>{\raggedright\arraybackslash}p{0.18\textwidth}
                  >{\raggedright\arraybackslash}p{0.25\textwidth}
                  >{\raggedright\arraybackslash}p{0.36\textwidth}
                  >{\raggedright\arraybackslash}p{0.13\textwidth}}
\toprule
Topic & Literature & Formalization & Classification \\
\midrule
\endhead
Continuity assumptions & Assumptions A2 and A4 are pointwise in the perturbed
coordinate. & The proved transfer theorem packages observations as
embedding--label pairs and uses joint continuity of the learning rule and loss.
Both the paper-shaped predicates and the stronger predicates used by the proof
are exposed. & Stronger sufficient assumptions \\
\addlinespace
Compactness assumption A3 & The paper states closedness, boundedness, and
completeness of the image. & The main proof uses a uniform compact-range
condition, equivalent in finite-dimensional Euclidean space; a literal
paper-shaped predicate is also retained. & Equivalent finite-dimensional
encoding \\
\addlinespace
Measurability & Mostly implicit. & Measurability of estimated embeddings is
explicit because risks and expectations are taken. & Surfaced technical
assumption \\
\addlinespace
Maximum and budgets & Equation (3) uses a finite maximum and two asymptotic
budgets. & The formalization uses a finite supremum and one abstract budget
index, which may encode a joint schedule. & Representation change \\
\addlinespace
MDS estimator used by the bridge & The argument cites the eigengap-free
asymptotic raw-stress consistency result from the 2024 paper. & The current
end-to-end bridge realizes the practical estimator as classical spectral MDS
and therefore requires a uniform lower population eigenvalue bound.  This
assumption is not stated in the Helm paper. & Material surfaced discrepancy \\
\bottomrule
\end{longtable}

The last row is the most important discrepancy in this paper.  There are two
honest ways to close it: connect Helm directly to the raw-stress consistency
route used in the paper, or amend the spectral-estimator theorem with the
necessary eigenvalue-stability condition.  The current formalization takes the
second route and states the condition explicitly.


\paragraph{Role in the Perfect Quench scaffold.}
No new inference-transfer theorem is required for the fixed-subset Quench
result.  The reusable contribution from this paper track is its discipline
around continuity, compactness, and measurable event composition.  The ledger
therefore records no artificial Helm-specific dependency; future reuse should
be added only when it removes a Quench-facing hypothesis.

\section{Helm, Johnson, and Priebe (2026): Quench}

\paragraph{Source asset policy.}
The repository retains both the paper's exact Markdown transcription and its
original \LaTeX{} source.  A second standalone distilled note is not currently
useful: Theorem~2 and its proof are short, contiguous, and already preserved in
source order.  This section is therefore the transformative mathematical
reconstruction.  It modernizes the argument, identifies hidden steps, and maps
them to the preferred \Lean{} declarations without duplicating the paper's
prose.

\begin{longtable}{>{\raggedright\arraybackslash}p{0.18\textwidth}
                  >{\raggedright\arraybackslash}p{0.25\textwidth}
                  >{\raggedright\arraybackslash}p{0.36\textwidth}
                  >{\raggedright\arraybackslash}p{0.13\textwidth}}
\toprule
Topic & Literature & Formalization & Classification \\
\midrule
\endhead
Estimator and ties & The displayed estimator averages all references tied for
minimum estimated distance, but the proof immediately replaces it by one
chosen nearest reference. & The preferred theorem proves the bound for the
literal tie average.  Every tied reference satisfies the same distance bound,
so convex averaging preserves the pointwise score estimate. & Faithfulness
repair \\
\addlinespace
Theorem quantifiers & Theorem~2 says that for every error tolerance there exists
a triple $(n,m,r)$ giving small MSE, then interprets this as eventual
query-efficiency. & The formal result fixes a proper query subset and supplies
explicit stage-dependent sample schedules.  It proves an eventual
high-probability risk comparison, then lifts this statement to every proper
subset with the same quantifier order used by the paper's definition. &
Stronger and more explicit conclusion \\
\addlinespace
Pointwise error to MSE & The proof obtains a high-probability pointwise error
bound and then states the corresponding MSE bound.  Uniformity in the target
model and measurability of the good event are implicit. & The formal proof
constructs a measurable event on which the squared error bound holds for every
target model, then integrates that uniform bound against $P_f$. & Missing step
made explicit \\
\addlinespace
Assumption~1 constant & The paper states existence of some $\gamma>0$.  Earlier
interfaces additionally asked callers to prove positivity of the particular
stored constant. & The public certificate now accepts any constant satisfying
the displayed Lipschitz inequality and internally replaces it by
$\max\{\gamma,1\}$.  The separate positivity field is gone from both the
paper-shaped and Perfect Quench interfaces. & Redundant hypothesis removed \\
\addlinespace
Assumption~2 & The paper says that every compact subset has positive measure and
calls this equivalent to positive probability of every target-centered ball.
The first statement is not literally correct, since it includes the empty set
and may include zero-mass compact sets. & \code{PerspectiveFullSupport} records
exactly the property used by the proof: every positive-radius perspective ball
centered at a model has positive $P_f$ mass.  Compactness plus finite covering
then upgrades pointwise positive mass to uniform high-probability coverage. &
Measure-theoretic repair \\
\addlinespace
Reference sampling & Iid reference sampling is prose-level. &
\code{IIDReferenceSampler} records the finite joint-product law.  Coverage is a
theorem derived from iid sampling and full support rather than a second
probabilistic assumption. & Surfaced probabilistic interface \\
\addlinespace
Reference and response randomness & The source proof does not separate the
random reference sample from cached-response noise. & Perfect Quench uses a
product probability space.  This makes independence structural and prevents a
randomly selected reference model from biasing the response array selected for
averaging. & Model clarification \\
\addlinespace
Population perspective geometry & The paper treats the true perspective as the
population target of the estimated DKPS coordinates. & The formal theorem asks
for one model-level response-distance realization.  Centering, double
centering, the CMDS Gram identity, positive semidefiniteness, rank, and the
reference-to-target radial identity are derived internally. & Hidden bridge
made explicit; downstream hypotheses removed \\
\addlinespace
Spectral regularity & The concentration theorem is invoked ``under technical
assumptions,'' with the stable rank-$d$ geometry left implicit. & One population
covariance floor, \code{PerspectiveNondegeneracy}, yields measurable
high-probability lower and upper spectral events.  No eigenvalue bound is
assumed for every reference outcome. & Technical assumption isolated and
samplewise assumptions removed \\
\addlinespace
Response concentration & The paper imports a representation-concentration
bound and uses the qualitative regime $r=\omega(n^3)$. & Perfect Quench starts
from measurable pairwise-independent raw replicates with a uniform second
moment, constructs their averages, and proves finite-class or entropy-net
uniform concentration using conservative explicit schedules. & Stronger input
model and constructive rate proof \\
\addlinespace
Finite versus infinite model classes & The paper speaks about a general model
population without specifying how uniform control over all target models is
obtained. & There are two complete capstones: a finite-model theorem using a
finite union bound, and a compact-infinite theorem using pathwise response
regularity and polynomial Euclidean covers. & Ambiguous scope split into two
honest theorems \\
\addlinespace
All query subsets & Query efficiency ranges over all budgets below the full
benchmark. & The fixed-subset theorem is lifted to size-$m$ and all-proper-
subset predicates.  Each subset may have its own perspective and response
certificate, matching the source quantifier order; no unjustified uncountable
or combinatorial event intersection is introduced. & Quantifier-faithful
completion \\
\addlinespace
Baseline positivity & The paper explicitly restricts the comparison to
$\operatorname{MSE}(\hat y_Q)>0$. & The same condition is retained.  It is
mathematically necessary for deriving strict eventual improvement solely from
convergence of the new estimator's MSE to zero. & Faithful necessary hypothesis
\\
\bottomrule
\end{longtable}

\section{Perfect Quench theorem family: completed state}

The \code{DkpsQuench/Perfect} development is no longer a scaffold.  Every proof
body in the directory is closed, including the finite, compact-infinite, and
all-proper-subset capstones:
\begin{itemize}
\item \code{perfectQuench\_finite\_fixedSubset};
\item \code{perfectQuench\_infinite\_fixedSubset};
\item \code{perfectQuench\_finite\_allQueries};
\item \code{perfectQuench\_infinite\_allQueries}.
\end{itemize}

\begin{longtable}{>{\raggedright\arraybackslash}p{0.22\textwidth}
                  >{\raggedright\arraybackslash}p{0.31\textwidth}
                  >{\raggedright\arraybackslash}p{0.37\textwidth}}
\toprule
Track & Proved mathematical content & Caller-visible data eliminated \\
\midrule
\endhead
Population geometry & Exact finite centering, Euclidean double centering,
configuration Gram realization, PSD/rank, and target radial identities. &
Separate configuration, Gram, radial, PSD, and rank witnesses collapse to one
model-level response realization. \\
\addlinespace
Raw responses & Product-space iid lifting, concrete replicate means,
second-moment transfer through random reference selection, finite union bounds,
and finite-net response events. & Abstract response means, hand-built moment
events, and per-coordinate augmented measurability certificates. \\
\addlinespace
Spectral regularity & Population covariance algebra, coordinate weak laws,
entrywise-to-quadratic transfer, centered scatter identities, rectangular
scatter--Gram spectral transport, target augmentation, and measurable spectral
subevents. & Global lower and upper eigenvalue bounds for every sample outcome.
\\
\addlinespace
Compactness and infinite classes & Response envelopes, raw-to-sample and
raw-to-population Lipschitz transfer, polynomial covers, entropy union bounds,
and deterministic extension from the stage net. & Caller-selected nets,
covering constants, entropy exponents, response envelopes, and separate sample
and population regularity certificates. \\
\addlinespace
Rates and composition & Explicit safe tolerances, finite and entropy-aware
replicate schedules, event intersections, CMDS perturbation, tie averaging, and
risk conversion. & Entry-rate arithmetic, polar smallness bookkeeping, and the
full internal growing-configuration certificate. \\
\addlinespace
Query budgets & Fixed-subset results and direct lifts to every proper subset. &
No common event across all subsets is demanded; the source's subset-specific
eventual thresholds are preserved. \\
\bottomrule
\end{longtable}

The safe schedules are deliberately existence-oriented.  They are explicit and
proved sufficient, but they are not advertised as a sharp restatement of the
paper's $r=\omega(n^3)$ regime.  A future rate theorem should parameterize the
capstone by abstract asymptotic conditions and recover both the safe schedule
and sharper literature schedules as corollaries, without changing the
probabilistic or geometric core.

\section{Current Quench hypothesis audit}

The preferred public interfaces now have the smallest hypothesis load that the
current proof architecture justifies.  The latest reduction removes the
standalone positivity proof for the stored score-Lipschitz constant: any valid
constant is enlarged internally to $\max\{\gamma,1\}$ before the quantitative
nearest-neighbor argument divides by it.

The remaining assumptions divide as follows.

\begin{longtable}{>{\raggedright\arraybackslash}p{0.23\textwidth}
                  >{\raggedright\arraybackslash}p{0.30\textwidth}
                  >{\raggedright\arraybackslash}p{0.37\textwidth}}
\toprule
Hypothesis & Present status & Reduction assessment \\
\midrule
\endhead
Measurable perspective & Used by perspective-ball events, covariance sampling,
and pushforward/integration arguments. & Retain for now.  It is not a formal
consequence of pairwise distance realization alone. \\
\addlinespace
Perspective full support & Exactly supplies nearest-reference coverage. &
Retain unless the conclusion is weakened to targets in the support or coverage
is supplied directly. \\
\addlinespace
Iid reference sampler & Supplies the product miss-probability formula. & Retain
for the current theorem.  Mixing or exchangeable samplers require a separate
coverage theorem, not deletion of the sampling hypothesis. \\
\addlinespace
Raw response measurability, $L^2$ control, pairwise independence, and second
moment & These are the concrete assumptions from which the response-mean event
is proved. & Retain the mathematical content.  Pairwise independence is already
weaker than full independence. \\
\addlinespace
Measurability of the supplied population mean & Stored in
\code{RawIIDResponseModel}. & Candidate for removal.  Joint measurability of the
raw response and the coordinate mean identities should permit a parameter-
integral measurability proof using the product-integral machinery already used
in \code{RawResponses.lean}. \\
\addlinespace
Exact response-distance realization & Connects the population response metric
to the supplied perspective. & Conceptually necessary unless the perspective
is defined canonically from the population responses.  The next architectural
improvement is a canonical constructor that turns realization from an
assumption into a theorem. \\
\addlinespace
Population covariance floor & Provides identifiability of all retained
perspective directions and the stable rank-$d$ CMDS geometry. & Cannot be
removed without weakening the conclusion.  Existing finite-frame tools can,
however, derive it from concrete spanning/frame assumptions in applications.
\\
\addlinespace
Centered-mean field in the nondegeneracy certificate & The structure
currently records that each coordinate centered at \code{perspectiveMean} has
zero integral. & Candidate for removal.  The covariance-floor development
already derives integrability of every centered squared linear form; $L^2$ to
$L^1$ plus the definition of the Bochner mean should prove the field
automatically. \\
\addlinespace
Compact perspective range and pathwise raw-response Lipschitzness & Used only by
the compact-infinite route to construct polynomial nets and extend finite-net
concentration. & Retain for this route.  Alternate entropy or stochastic-
equicontinuity interfaces can be added as parallel capstones rather than
weakening these hypotheses syntactically. \\
\addlinespace
Positive baseline MSE & Converts convergence to zero into strict improvement
over the subset-score baseline. & Retain; this is explicit in the paper and
logically necessary for the stated comparison. \\
\addlinespace
Positive number of queried response rows & Needed because the response metric
uses the inverse of the query count. & Retain or redefine the zero-query case;
it is not an artifact of the proof. \\
\bottomrule
\end{longtable}

The two most promising next reductions using machinery already present in the
repository are therefore: (1) derive population-mean measurability from the
jointly measurable raw response family, and (2) derive the centered-mean
identity inside \code{PerspectiveNondegeneracy}.  Neither should be replaced by
a stronger application-facing assumption.

\section{Audit checklist for future edits}

For every new preferred capstone, record:
\begin{enumerate}
\item the exact paper theorem or algorithm step it represents;
\item each assumption not written in the source paper;
\item each paper assumption discharged internally by a preceding theorem;
\item whether the estimator is raw-stress MDS, classical MDS, or another
      explicitly named variant;
\item whether the conclusion is fixed-target, uniform over targets, fixed-query,
      or uniform over query subsets;
\item whether rates are qualitative, explicit but loose, or quantitatively
      matched to the paper;
\item any finite-dimensional, finite-model, measurability, separability, or
      compactness restriction.
\end{enumerate}

\begin{thebibliography}{9}
\bibitem{acharyya2024}
A. Acharyya, M. W. Trosset, C. E. Priebe, and H. S. Helm.
\newblock Consistent estimation of generative model representations in the data
kernel perspective space.
\newblock arXiv:2409.17308, 2024.

\bibitem{acharyya2025}
A. Acharyya, J. Agterberg, Y. Park, and C. E. Priebe.
\newblock Concentration bounds on response-based vector embeddings of black-box
generative models.
\newblock arXiv:2511.08307, 2025.

\bibitem{helm2025}
H. S. Helm, A. Acharyya, Y. Park, B. Duderstadt, and C. E. Priebe.
\newblock Statistical inference on black-box generative models in the data
kernel perspective space.
\newblock Findings of ACL, 2025.

\bibitem{quench2026}
H. Helm, B. Johnson, and C. E. Priebe.
\newblock Query-efficient model evaluation using cached responses.
\newblock arXiv:2605.07096, 2026.
\end{thebibliography}

\end{document}
